\documentclass[11pt]{article}
\usepackage[margin=1in]{geometry}
\usepackage{amsmath,amssymb,amsthm,mathtools,graphicx,booktabs}
\usepackage[hidelinks]{hyperref}
\newtheorem{theorem}{Theorem}
\newtheorem{lemma}{Lemma}
\newtheorem{remark}{Remark}
\newcommand{\T}{\mathbb T}
\newcommand{\M}{\mathcal M}
\newcommand{\F}{\mathcal F}
\newcommand{\E}{\mathbb E}
\newcommand{\tr}{\operatorname{Tr}}
\newcommand{\BMO}{\operatorname{BMO}}
\setlength{\parindent}{0pt}
\setlength{\parskip}{0.62em}

\title{Failure of Unconditional Extrapolation for\\Operator-Valued Dyadic Paraproducts}
\author{Run \texttt{fa\_banach\_001}, lane 10}
\date{11 August 2026}

\begin{document}
\maketitle

\textbf{Status.} Candidate full counterexample, likely valid, for expert review.

\section{Source question}

Let $(\F_j)_{j\geq0}$ be the usual dyadic filtration on the circle and
write $E_j$ and $d_j=E_j-E_{j-1}$ for conditional expectation and
martingale difference.  For a von Neumann algebra $\M$ with a normal
semifinite faithful trace, Mei defines the left dyadic paraproduct
\[
 \pi_\varphi f=\sum_j(d_j\varphi)(E_{j-1}f)
 \quad\text{on }L^p(\T;L^p(\M)).
\]
Theorem~1.2 of \cite{Mei} proves that boundedness for one
$1<p<\infty$ extrapolates to all such $p$ provided
$\varphi\in\BMO_{\M}(\T,\M)$.  The paper then asks whether that
operator-norm BMO assumption can be removed
\cite[Open Question, PDF p.~9]{Mei}.

\begin{center}
 \includegraphics[width=\textwidth]{figures/open_problem_crop.png}
\end{center}

Equivalently, the question asks whether boundedness of $\pi_\varphi$ on
one $L^p(\T;L^p(\M))$ alone forces boundedness on every such space.

\section{Claimed contribution}

\begin{theorem}\label{thm:main}
There are a semifinite tracial von Neumann algebra $\M$ and an
$\M$-valued Bochner integrable symbol $\varphi$ such that
\[
 \|\pi_\varphi\|_{L^2(\T;L^2(\M))\to L^2(\T;L^2(\M))}\leq1,
\]
but, for every $2<p<\infty$,
\[
 \pi_\varphi:L^p(\T;L^p(\M))\longrightarrow
 L^p(\T;L^p(\M))
 \quad\text{is unbounded}.
\]
The symbol belongs to $L^q(\T;\M)$ in operator norm for every finite
$q$, but $\|\varphi\|_{\BMO_{\M}}=\infty$.  Hence the open question has
a negative answer.
\end{theorem}

\section{Proof intuition}

The one-sided paraproduct can see columns and rows very differently.  In an
$n\times n$ matrix corner, put successive martingale differences in the
first row: $d_k\varphi=r_k e_{1k}$.  On $L^2$, martingale orthogonality and
the identity $\|e_{1k}x\|_2=\|e_{kk}x\|_2$ turn the output norm into a sum
of mutually orthogonal input rows, so the norm is at most one.  On the
identity matrix, however, the output is the row
$(r_1,\ldots,r_n)$, whose only singular value is $\sqrt n$.  Its ratio to
the Schatten-$p$ norm of the identity is $n^{1/2-1/p}$.

A single finite block gives only a large norm, not an unbounded operator.
We place the $n$th block on its own exponentially short dyadic interval and
use mutually orthogonal matrix corners.  The resulting infinite symbol is
Bochner integrable, retains the uniform $L^2$ contraction, and contains all
the divergent $L^p$ tests simultaneously.

\section{Construction}

Identify $\T$ with $[0,1)$ and let
\[
 \mathcal H=\bigoplus_{n\geq1}\mathbb C^n,
 \qquad \M=B(\mathcal H),
\]
equipped with its usual trace.  In the $n$th summand let
$e_{ab}^{(n)}$ denote the standard matrix units and let
$P_n=\sum_{k=1}^n e_{kk}^{(n)}$.

For $n\geq1$, set
\[
 I_n=[2^{-n},3\cdot2^{-(n+1)}),\qquad m_n=n+1.
\]
Thus $I_n$ is an atom of $\F_{m_n}$, has length $2^{-(n+1)}$, and the
intervals $I_n$ are pairwise disjoint.  On $I_n$ let
$r_{n,1},\ldots,r_{n,n}$ be the local Rademacher functions: after extension
by zero outside $I_n$, they satisfy
\[
 r_{n,k}\text{ is }\F_{m_n+k}\text{-measurable},\qquad
 E_{m_n+k-1}r_{n,k}=0,\qquad |r_{n,k}|=1_{I_n}.
\]
Define
\begin{equation}\label{eq:symbol}
 \varphi(t)=\sum_{n=1}^{\infty}\sum_{k=1}^n
 r_{n,k}(t)e_{1k}^{(n)}.
\end{equation}
At any $t$, at most one outer summand is nonzero.  On $I_n$ the operator
$\varphi(t)$ is a rank-one row of length $\sqrt n$, so
\[
 \|\varphi(t)\|_{\M}=\sqrt n\quad(t\in I_n).
\]
Consequently
\[
 \int_\T\|\varphi(t)\|_{\M}^q\,dt
 =\sum_{n\geq1}2^{-(n+1)}n^{q/2}<\infty
 \qquad(1\leq q<\infty).
\]
In particular, \eqref{eq:symbol} is a legitimate Bochner integrable
$\M$-valued symbol.

The placement in the filtration gives the exact martingale differences
\begin{equation}\label{eq:diff}
 d_j\varphi
 =\sum_{\substack{n\geq1,\ 1\leq k\leq n\\m_n+k=j}}
 r_{n,k}e_{1k}^{(n)}.
\end{equation}
The sum in \eqref{eq:diff} is finite for each $j$.  Indeed, coarser local
Rademachers are already $\F_{j-1}$-measurable, while finer ones have zero
conditional expectation at both levels $j-1$ and $j$.

\section{The uniform \texorpdfstring{$L^2$}{L2} bound}

\begin{lemma}\label{lem:l2}
The series defining $\pi_\varphi$ converges on
$L^2(\T;S^2(\mathcal H))$ and $\|\pi_\varphi\|_{2\to2}\leq1$.
\end{lemma}

\begin{proof}
Start with a finite martingale $f$.  Terms belonging to different global
martingale levels are orthogonal in $L^2$.  At a common level, terms from
different $n$ have orthogonal left ranges.  Hence \eqref{eq:diff} gives
\begin{align*}
 \|\pi_\varphi f\|_2^2
 &=\sum_{n\geq1}\sum_{k=1}^n
   \bigl\|r_{n,k}e_{1k}^{(n)}E_{m_n+k-1}f\bigr\|_2^2. \tag{1}
\end{align*}
All norms here include integration in $t$ and the Hilbert--Schmidt norm.
Because $I_n\in\F_{m_n}$ and $m_n+k-1\geq m_n$,
\[
 1_{I_n}E_{m_n+k-1}f
 =E_{m_n+k-1}(1_{I_n}f).
\]
For every Hilbert--Schmidt operator $x$,
\[
 \|e_{1k}^{(n)}x\|_2=\|e_{kk}^{(n)}x\|_2.
\]
Conditional expectation is an $L^2$ contraction and commutes with the
constant matrix unit, so each summand in (1) is at most
\[
 \|1_{I_n}e_{kk}^{(n)}f\|_2^2.
\]
The intervals $I_n$ are disjoint and all the projections
$e_{kk}^{(n)}$ are mutually orthogonal.  Therefore
\[
 \|\pi_\varphi f\|_2^2
 \leq\sum_{n\geq1}\sum_{k=1}^n
 \|1_{I_n}e_{kk}^{(n)}f\|_2^2
 \leq\|f\|_2^2.
\]
The same estimate for partial sums makes the full orthogonal series
convergent.  Density then gives the asserted contraction on all of $L^2$.
\end{proof}

\section{Failure on every \texorpdfstring{$L^p$}{Lp} above two}

Fix $n$ and use the $\F_{m_n}$-measurable test function
\[
 f^{(n)}(t)=1_{I_n}(t)P_n.
\]
For every $1\leq k\leq n$,
$E_{m_n+k-1}f^{(n)}=f^{(n)}$.  Thus
\begin{equation}\label{eq:output}
 \pi_\varphi f^{(n)}
 =\sum_{k=1}^n r_{n,k}e_{1k}^{(n)}.
\end{equation}
The projection $P_n$ has $n$ singular values equal to one, whereas the
row in \eqref{eq:output} has exactly one nonzero singular value, namely
$\sqrt n$.  It follows that, for every finite $p$,
\[
 \|f^{(n)}\|_{L^p(\T;S^p)}=|I_n|^{1/p}n^{1/p},
 \qquad
 \|\pi_\varphi f^{(n)}\|_{L^p(\T;S^p)}
 =|I_n|^{1/p}\sqrt n.
\]
Hence
\begin{equation}\label{eq:ratio}
 \frac{\|\pi_\varphi f^{(n)}\|_p}{\|f^{(n)}\|_p}
 =n^{1/2-1/p}.
\end{equation}
For each fixed $p>2$, the right-hand side tends to infinity.  Therefore
$\pi_\varphi$ is unbounded on every $L^p(\T;S^p)$ with $p>2$.
Together with Lemma~\ref{lem:l2}, this proves Theorem~\ref{thm:main}.

Finally, the construction lies exactly outside the hypothesis in the source.
Since every local Rademacher has mean zero on $I_n$, one has
$\varphi_{I_n}=0$, and therefore
\[
 \left(\frac1{|I_n|}\int_{I_n}
       \|\varphi-\varphi_{I_n}\|_{\M}^2\,dt\right)^{1/2}
 =\sqrt n.
\]
Taking the supremum over dyadic intervals proves
$\|\varphi\|_{\BMO_{\M}}=\infty$.

\begin{remark}[Finite tracial variant]
If a finite trace is preferred, replace $B(\mathcal H)$ by
$\prod_{n\geq1}M_n$ with
$\tau((x_n))=\sum_{n\geq1}2^{-n}\operatorname{tr}_n(x_n)$, where
$\operatorname{tr}_n$ is normalized.  The same symbol and the same proof
work blockwise.  Both norms in the $n$th test acquire the same factor
$2^{-n/p}$, so the ratio \eqref{eq:ratio} is unchanged.
\end{remark}

\section{Verification, limitations, and novelty status}

The script \texttt{code/verify\_counterexample.py} independently checks the
local martingale-difference identities, random finite-block versions of the
$L^2$ contraction for $n\leq8$, the rank-one singular-value computation, and
the exact ratios \eqref{eq:ratio} for several $p$.  These computations are
sanity checks; the proof above is exact and does not depend on them.

The construction answers only the yes/no extrapolation question.  It does not
characterize all symbols whose paraproducts extrapolate, nor does it settle
the still broader problem of characterizing $L^2$ boundedness of
operator-valued paraproducts.

A bounded novelty search was performed on 11 August 2026: the run's title,
abstract, source-signal, result, and attempt indexes were searched for the
arXiv id and the core phrases; exact and keyword searches were made for an
operator-valued paraproduct bounded on $L^2$ but unbounded on $L^p$; and
later arXiv papers citing the source were inspected.  In particular,
Wei--Zhang still state that general $L^2$ boundedness characterization is
open and formulate extrapolation with the same strong BMO hypothesis
\cite{WeiZhang}.  No matching counterexample was located.  This is bounded
negative evidence, not a priority claim.

\textbf{Human-review recommendation.} Send to a specialist in
noncommutative martingales.  The critical audit points are the global
martingale-difference identity \eqref{eq:diff}, the orthogonality in
Lemma~\ref{lem:l2}, and whether the source question is read in its literal
ambient class of Bochner integrable $\M$-valued symbols.

\begin{thebibliography}{9}
\bibitem{Mei}
T.~Mei,
\emph{An extrapolation of operator valued dyadic paraproducts},
J. London Math. Soc. (2) 81 (2010), no.~3, 650--662;
\href{https://arxiv.org/abs/0709.4229}{arXiv:0709.4229}.

\bibitem{WeiZhang}
Z.~Wei and H.~Zhang,
\emph{Boundedness of operator-valued commutators involving martingale
paraproducts},
\href{https://arxiv.org/abs/2401.08729}{arXiv:2401.08729}.
\end{thebibliography}

\end{document}
